\chapter{洁净性}\label{ch:6}\index{洁净性|hyperpage}
% 原PDF第36页中部
\section{开子叠 \texorpdfstring{$C_k$}{Cₖ} 上的 Fourier 塔}\label{sec:6.1}
本章展开 Fourier 变换递归。对 $1\le k\le n$，令
\[\Coh_k:=\{M_k\in\Coh(X)\mid\rank(M_k)=k\}.\]
有两个截面叠：
\[\Coh'_k:=\{\Omega_X^{k-1}\to M_k\text{ 在一般点单射}\},\qquad\overline{\Coh}'_k:=\{\Omega_X^{k-1}\to M_k\}.\]
上横线允许截面不是单射或为零。固定开子叠 $C_k\subset\Coh_k$，并记
\[C'_k:=\Coh'_k\times_{\Coh_k}C_k,\qquad\overline C'_k:=\overline{\Coh}'_k\times_{\Coh_k}C_k.\]
要求这些开子叠满足以下条件。
\begin{enumerate}
\item\label{cond:6.1.1} 对 $M_k\in C_k$，有 $\Ext^1(\Omega_X^{k-1},M_k)=0$。因此，$\overline C'_k\to C_k$ 是纤维为 $\Hom(\Omega_X^{k-1},M_k)$ 的向量丛。
\item\label{cond:6.1.2} 若 $M_k^{\mathrm{tf}}$ 是无挠商，则
\[\ndeg(M_k^{\mathrm{tf}})>nk(2g-2).\]
证明使用这一归一化次数形式。
\item\label{cond:6.1.3} 若 $M_k\in C_k$ 且 $\Omega_X^{k-1}\hookrightarrow M_k$ 在一般点单射，则 $M_k/\Omega_X^{k-1}\in C_{k-1}$。
\end{enumerate}
由 Serre 对偶，$\overline C'_k\to C_k$ 的线性对偶的纤维为
\[\Ext^1(M_k,\Omega_X^k).\]
记这一对偶向量丛为 $\mathcal E_{k+1}\to C_k$；它的点是扩张类
\[0\longrightarrow\Omega_X^k\longrightarrow M_{k+1}\longrightarrow M_k\longrightarrow0.\]
这就是下一次 Fourier 变换从层级 $k$ 移到层级 $k+1$ 的原因。

归纳定义 $F_k$。首先，
\[F_1:=q_1^*\cL_E[\text{移位}](\text{扭曲}),\]
其中 $q_1:C'_1\to\Tor$ 取 $\cO_X=\Omega_X^0\hookrightarrow M_1$ 的商。若 $j_k:C'_k\hookrightarrow\overline C'_k$ 为开嵌入，则
\[F_{k+1}:=\mathsf F_{\overline C'_k/C_k}(j_{k,!}F_k)|_{\mathcal E_{k+1}^{\ne0}}.\]
这里 $\mathcal E_{k+1}^{\ne0}$ 是中间项属于 $C_{k+1}$ 且典范截面 $\Omega_X^k\to M_{k+1}$ 在一般点单射的开部分。不会造成混淆时，借助从扩张类到中间项的遗忘映射，将这一限制看作 $C'_{k+1}$ 上的层。函子 $\mathsf F_{\overline C'_k/C_k}$ 是 $C_k$ 上两个对偶向量丛之间的 Fourier–Deligne 变换。
\begin{transnote}{含糊点拨}
这里的「相对」有两层：$\mathcal E_{k+1}$ 与 $\overline C'_k$ 都是 $C_k$ 上的向量丛（后者是其紧化），所以变换是逐纤维做的 §5.2，移位与扭转按相对维数取，而 $C_k$ 本身不动。先取 $j_{k,!}$ 再变换，是为了让结果只落在 $\mathcal E_{k+1}^{\ne 0}$ 上：中间项在一般点单射正是下一层叠 $C'_{k+1}$ 的定义条件，缺了 $!$-延拓，支撑会伸到有挠缺陷的边界上。[依据：原书 p.37–38 本段与 §6.2 的定义；并见 [KL85] 的 Fourier 变换。]
\end{transnote}
\srcpage{37}

在 $\Bun'_n\cap C'_n$ 上，结果是初步的自守层：
\[F_n|_{\Bun'_n\cap C'_n}=\Aut'_E|_{\Bun'_n\cap C'_n}.\]
更一般地，$F_n$ 由紧化 Whittaker 叠计算：
\[Q_n=\left\{\begin{array}{l}0=M_0\subset M_1\subset\cdots\subset M_n,\quad M_i/M_{i-1}\simeq\Omega_X^{n-i},\\M_n\hookrightarrow M\in C_n\text{，余核为挠层}\end{array}\right\}.\]
它取代未紧化的 Whittaker 叠。由于必须保留挠缺陷，允许目标属于 $\Coh_n$，而不仅属于 $\Bun_n$。
\section{洁净性定理}\label{sec:6.2}
在形成下一次 Fourier 变换所需的层级上，目标为：
\begin{theorem}{洁净性}{6.2.1}
对 $1\le k<n$，$F_k$ 跨越非单射轨迹的延拓是洁净的，即在 $\overline C'_k$ 上有
\[j_{k,!}F_k\simeq j_{k,*}F_k.\]
因此，Fourier 归纳在所选开子叠上保持所需的不可约反常层。
\end{theorem}

按挠部分的长度过滤。置
\[\Coh_k^{\mathrm{tor},\le\ell}:=\{M_k\in\Coh_k\mid\length(\operatorname{Tor}(M_k))\le\ell\},\qquad C_k^{\mathrm{tor},\le\ell}:=C_k\cap\Coh_k^{\mathrm{tor},\le\ell}.\]
则 $\Coh_k^{\mathrm{tor}=0}=\Bun_k$，$C_k=\bigcup_{\ell\ge0}C_k^{\mathrm{tor},\le\ell}$。对 $\ell$ 归纳，证明 $\overline C'_k{}^{\mathrm{tor},\le\ell}$ 上的洁净性。
\section{基础情形：无挠层}\label{sec:6.3}
\begin{proof}
取 $\ell=0$。从线丛到向量丛的映射要么为零，要么为单射。因此，$C'_k{}^{\mathrm{tor}=0}$ 在 $\overline C'_k{}^{\mathrm{tor}=0}$ 中的补集是零截面
\[i:C_k^{\mathrm{tor}=0}\hookrightarrow\overline C'_k{}^{\mathrm{tor}=0}.\]
令 $q:\overline C'_k{}^{\mathrm{tor}=0}\to C_k^{\mathrm{tor}=0}$ 为向量丛投影。该层沿纤维上的伸缩作用的轨道局部常值。与 $\GL_2$ 情形一样，洁净性等价于边界消失 $i^*j_{k,*}F_k=0$；由 Verdier 对偶，只需对 $F_k$ 及其对偶证明紧支撑版本。收缩原理给出
\[i^!j_{k,!}F_k\simeq q_!j_{k,!}F_k.\]
\begin{transnote}{易错预警}
用 $i^!$ 而不是 $i^*$：$i$ 是零截面的闭嵌入，$i^!$ 才天然配对紧支撑推前，收缩原理给出的也正是 $!$-版本。经 Verdier 对偶，$i^!j_{k,!}\simeq q_!j_{k,!}$ 对偶到 $i^*j_{k,*}\simeq q_*j_{k,*}$，所以只需证 $i^*j_{k,*}F_k=0$ 就能同时得到 $j_!\simeq j_*$，即 $j_{!*}=j_!=j_*$。前提是 $F_k$ 沿 $\Gm$-伸缩轨道局部常值（monodromic），否则收缩原理不适用。[来源：nLab "vanishing cycle"，https://ncatlab.org/nlab/show/vanishing+cycle（Last revised 2015-01-25；2026-11-17 访问）；并见原书 p.39 本段。]
\end{transnote}
所以只需证明 $q_!j_{k,!}F_k=0$；将同样的论证用于对偶，得到 $*$ 边界消失。

使用 $F_k$ 的 Whittaker 公式，这个对象为
\[q_!j_{k,!}F_k\simeq(r\circ\pi)_!(p_1^*\cL_E\otimes p_2^*\cL_\psi)[\text{移位}](\text{扭曲})|_{C_k^{\mathrm{tor}=0}},\]
其中 $p_1$ 记录挠商，$p_2$ 为 Whittaker 特征标。在无挠零截面边界上，紧化 Whittaker 叠限制为普通 Whittaker 叠
\[\Bun_{N_k}\xrightarrow{\pi}\Bun'_k\xrightarrow{r}\Bun_k.\]
令 $W_k$ 为基础 Whittaker 层，即将 $\exp_\psi$ 从 $\Bun_{N_k}$ 直接像到 $\Bun_k$ 所得的层，采用相同的移位与扭曲。在归一化次数 $d$ 的分支上，
\begin{transnote}{含糊点拨}
$\exp_\psi$ 指把 $\Ga$ 上的 Artin–Schreier 层沿 $\Bun_{N_k}$ 的简单根坐标拉回所得的层，其直像即 Whittaker 层；这里取 $*$-直像而非 $!$-直像，是因为 $W_k$ 要作为非紧支撑的核参与后面的 $!$-卷积。移位与扭转跟 $\Bun_k$ 的反常 $t$-结构对齐，相对维数与 §5.2 中 $[r](r/2)$ 的 $r$ 同一个来源：$\Bun_{N_k}\to\Bun_k$ 的纤维维数。[依据：原书 p.39 本段；并见 p.22 §3.1 对 Whittaker 特征标的定义。]
\end{transnote}
\[q_!j_{k,!}F_k\simeq\Av_E^d(W_k)|_{C_k^{\mathrm{tor}=0}}.\]
\srcpage{38}
归一化次数的约定使 $\Bun_{N_k}$ 位于 $\Bun_k$ 的次数零分支之上。\hyperref[cond:6.1.2]{条件(2)}给出
\[C_k^{\mathrm{tor}=0}\subset\bigcup_{d>nk(2g-2)}\Bun_k^d.\]
由于所考虑的归纳层级满足 $k<n$，\thref{5.7.2}说明：对 $d>nk(2g-2)$，$\Av_E^d=0$。因此 $q_!j_{k,!}F_k=0$，基础情形得证。
\end{proof}
\section{归纳步骤与基变换}\label{sec:6.4}
\begin{proof}
假设已知 $C_k^{\mathrm{tor},\le\ell-1}$ 上的洁净性。为避免将目标层与 Whittaker 旗的终点混淆，将截面的目标记为 $M$。新的边界层为
\[Z:=\{\Omega_X^{k-1}\to M\mid\text{不是单射},\ \length(\operatorname{Tor}(M))=\ell\}.\]
等价地，在 $\overline C'_k{}^{\mathrm{tor},\le\ell}$ 中，已经处理的开部分为
\[U:=C'_k{}^{\mathrm{tor},\le\ell}\cup\overline C'_k{}^{\mathrm{tor},\le\ell-1},\]
$Z$ 为剩余的闭部分。

当 $\ell>0$ 时，不存在从挠长度不超过 $\ell$ 的轨迹到挠长度等于 $\ell$ 的轨迹的典范投影。置 $Y_\ell:=\overline C'_k{}^{\mathrm{tor},\le\ell}$，并通过基变换添加这样的表现。定义
\[\widetilde Y:=\left\{\begin{array}{l}\Omega_X^{k-1}\to M,\quad0\to\mathring M\to M\to\mathcal T\to0,\\\mathring M\in\Bun_k,\ \mathcal T\in\Tor^\ell,\ M\in C_k,\ \mathring M\oplus\mathcal T\in C_k\end{array}\right\}.\]
忘去表现的映射 $\widetilde Y\to Y_\ell$ 在挠长度 $\ell$ 的层上光滑。不带截面时，这就是通过向量丛的修改给出 $\Coh_k$ 中挠长度 $\ell$ 的层的通常表现；相对切空间为 $\Hom(\mathring M,\mathcal T)$，所以相对维数局部常值。

固定点轨迹为
\[\widetilde Z:=\{\Omega_X^{k-1}\to\mathcal T,\quad\mathring M\in\Bun_k,\ \mathring M\oplus\mathcal T\in C_k\}.\]
它通过分裂扩张 $M=\mathring M\oplus\mathcal T$ 映入 $\widetilde Y$。注意 $\widetilde Z\ne Z$：在 $\widetilde Z$ 中，分裂／表现是数据的一部分。有投影
\[\widetilde f:\widetilde Y\longrightarrow\widetilde Z,\]
\[(\Omega_X^{k-1}\to M,\ 0\to\mathring M\to M\to\mathcal T\to0)\longmapsto(\Omega_X^{k-1}\to M\to\mathcal T,\mathring M).\]
对 $\Ext^1(\mathcal T,\mathring M)$ 中的扩张类作标量乘法，同时对所提升截面的 $\mathring M$ 分量作标量乘法，定义 $\widetilde Y$ 上的 $\Gm$ 作用，并延伸为收缩到 $\widetilde Z$ 的 $\A^1$ 作用。只缩放扩张类，会在固定点留下一个不需要的截面的向量丛分量。因此，收缩原理将归纳步骤归约为
\[\widetilde f_!\widetilde j_!(F_k|_{\widetilde U})=0,\qquad\widetilde U:=U\times_{\overline C'_k}\widetilde Y.\]
\end{proof}
\begingroup\setlength{\intextsep}{3pt}\setlength{\abovedisplayskip}{3pt}\setlength{\belowdisplayskip}{3pt}\setlength{\belowdisplayshortskip}{2pt}
\titlespacing*{\section}{0pt}{1.5ex plus .5ex}{1ex plus .2ex}
\section{纤维计算与满射情形}\label{sec:6.5}
\begin{proof}
使用紧化 Whittaker 公式
\[F_k\simeq\pi_!(p_1^*\cL_E\otimes p_2^*\cL_\psi)[\text{移位}](\text{扭曲})|_{C'_k},\]
其中 $p_1$ 记录挠商，$p_2$ 是 Whittaker 特征标。基变换后，取定一点 $(s_{\mathcal T}:\Omega_X^{k-1}\to\mathcal T,\mathring M)\in\widetilde Z$。相关纤维由 Whittaker 旗以及固定扩张
\[0\longrightarrow\mathring M\longrightarrow M\longrightarrow\mathcal T\longrightarrow0\]
之上的映射 $M_k\to M$ 组成，
\srcpage{39}
使复合 $\Omega_X^{k-1}=M_1\to M_k\to M\to\mathcal T$ 为固定映射 $s_{\mathcal T}$。诱导映射 $M_k\to\mathcal T$ 未必为满射。

若 $M_k\to\mathcal T$ 为满射，令 $\mathring M_k:=\ker(M_k\to\mathcal T)$，则有正合图
\origdiagram{p039-exact.pdf}{0.4668}
纤维等价于普通修改叠
\[\{\mathring M_k\to\mathring M\mid\mathring M/\mathring M_k\in\Tor^d\}.\]
因此，所需 $!$ 推前是 $\Av_E^d$ 的一个茎。在归一化次数下，Whittaker 旗的终点 $M_k$ 满足 $\ndeg M_k=0$，故 $\ndeg\mathring M_k=-\ell$。由于 $\mathring M\oplus\mathcal T\in C_k$，\hyperref[cond:6.1.2]{条件(2)}给出 $\ndeg\mathring M>nk(2g-2)$。于是
\[d=\ndeg\mathring M-\ndeg\mathring M_k>nk(2g-2)+\ell.\]
所以 $d>nk(2g-2)$，消失定理给出 $\Av_E^d=0$。这处理了满射情形。
\end{proof}
\section{非满射情形的归约}\label{sec:6.6}
\begin{proof}
假设 $M_k\to\mathcal T$ 不是满射，令 $\mathcal T'\subsetneq\mathcal T$ 为其像。置 $\mathring M_k:=\ker(M_k\to\mathcal T')$。令 $M'\subset M$ 为 $\mathcal T'$ 在 $M\to\mathcal T$ 下的逆像。则数据经由下图分解：
\origdiagram{p039-three.pdf}{0.4713}
从而
\[0\longrightarrow M'/M_k\longrightarrow M/M_k\longrightarrow\mathcal T/\mathcal T'\longrightarrow0.\]

令 $H_{\mathcal T}$ 和 $H_{\mathcal T'}$ 为相应修改纤维。映射 $H_{\mathcal T}\to H_{\mathcal T'}$ 与挠层图相容：
\origdiagram{p039-torses.pdf}{0.3825}
其中 $\Tor^{\mathrm{ses}}$ 分类挠层的短正合列；$m$ 记录中间项，$(s,q)$ 记录子对象与商。在当前情形，它们分别是 $M/M_k$、$M'/M_k$ 与 $\mathcal T/\mathcal T'$。
\Needspace{3\baselineskip}

拉回到这个图后，$H_{\mathcal T}\to H_{\mathcal T'}$ 是仿射空间纤维化。仿射方向由诸如 $\Ext^1(\mathcal T/\mathcal T',\mathring M_k)$ 的扩张数据，以及在商 $\mathcal T'$ 上选择相容提升 $M_k\to M'$ 控制。因此，忽略通常的仿射丛移位
\srcpage{40}
与 Tate 扭转后，通过基变换，剩余项归约为沿剩余投影 $\mathrm{pr}$ 的推前：
\[\mathrm{pr}_!((s,q)_!m^*\cL_E|_{H_{\mathcal T'}})=0.\]
所需命题是在 $\Tor$ 所用商叠归一化中的因子化恒等式
\[(s,q)_!m^*\cL_E\simeq\cL_E\boxtimes\cL_E,\]
它将在\chref{7}证明。它将挠商分为 $M'/M_k$ 部分与剩余的 $\mathcal T/\mathcal T'$ 部分，从而将非满射情形归约为满射情形。在假定该命题的情况下，洁净性的证明完成。
\begin{transnote}{含糊点拨}
这条恒等式说明：沿挠商部分把 $\cL_E$ 积掉，与先把挠商拆成 $M'/M_k$ 和 $\mathcal T/\mathcal T'$ 两段、再分别积掉，结果相同——即 $\cL_E$ 关于 $E$ 的乘法是因子化的。证明推迟到 \chref{7}（$C_k$ 的因子化定义）：先把非满射情形归约为满射，满射时 $(s,q)$ 沿纤维平坦，$\cL_E$ 逐纤维常值，恒等式退化为 $E$ 的乘法性。[依据：原书 p.40 本段；并见 \chref{7} 的因子化恒等式。]
\end{transnote}
\end{proof}
\endgroup
